%!TEX root = main.tex

Our goal is to construct a connection mapping $cm'$ from $tmap$ and $ss$ (the sequence of statements in $M$). 
For this purpose, we define two relations $cm_1 \xrightarrow[tmap]{s} cm_2$ and $cm_1 \xrightarrow[tmap]{ss'} cm_2$, where $s$ is a statement and $ss'$ is a sequence of statements. Note that $cm_1 \xrightarrow[tmap]{s} cm_2$ and $cm_1 \xrightarrow[tmap]{ss'} cm_2$ are mutually recursive, since a \prettyll{when} statement involves statement sequences. 
Finally, let $cm'$ be the connection mapping satisfying ${\sf emp} \xrightarrow[tmap]{ss} cm'$, where ${\sf emp}$ is the empty connection mapping. 

In the sequel, we show how to define the relations $cm_1 \xrightarrow[tmap]{s} cm_2$ and $cm_1 \xrightarrow[tmap]{ss'} cm_2$. Similar to the construction of $tmap$, when defining the two relations, we should deal with the mismatch between the abstraction levels of connection mappings and $M$. 
%
Furthermore, we should take care of the flipping status of fields when defining the semantics of connection statements of the form \prettyll{r1 <= r2}, where \prettyll{r1, r2} are both references. For instance, in the following {\hifirrtl} program, the connect statement \prettyll{a <= b} will be split into the following two connections between ground-type components, \prettyll{b.f3 <= a.f1} and \prettyll{a.f2 <= b.f4}. Note that in \prettyll{b.f3 <= a.f1}, the connection flow is reversed since \prettyll{f1} is a flipped field. 
%
\begin{figure*}[htbp]
\begin{tabular}{c}
%[basicstyle=\footnotesize\ttfamily]
\begin{lstlisting}
    wire a: {flip f1: UInt<1>, f2: UInt<3>}
    wire b: {flip f3: UInt<1>, f4: UInt<3>}
    a <= b
\end{lstlisting}
\end{tabular}
\end{figure*}

The relations $cm_1 \xrightarrow[tmap]{s} cm_2$ and $cm_1 \xrightarrow[tmap]{ss'} cm_2$ are defined by the rules in Fig.~\ref{fig:firrtl-sem}. 

%Given $tmap$, the semantics of a statement $s$ in $M$ is defined by a relation $ct \xrightarrow[tmap]{s} ct'$. %, where $s$ is seen as a modifier of module graphs. 
%Moreover, the relation can be easily extended to sequences of statements, that is, $ct \xrightarrow[tmap]{ss} ct'$, where $ss$ is a sequence of statements. 

\begin{figure}[htbp]
  \small
  \begin{gather*}
    \infer[skip]
    {\seqq{ cm \xrightarrow[tmap]{\mbox{skip}} cm}}
    {\seqq{-}}
    \quad
    \\[1ex]
    \infer[wire]
    {\seqq{cm_1 \xrightarrow[tmap]{\mbox{wire } x: t} cm_2}}
   {
    \begin{array}{c}
    id(x) = v \wedge tmap(v)= (ft, ori) \wedge \sizeofftype(ft) = sz \\
    \forall n.\ 0 \le n < sz \Rightarrow cm_1(v, n) = \svnone  \\
%    \forall n.\ n \ge size\_of\_ftype(ft) \Rightarrow vm(v, n) =  vm'(v,n)\\
%    \forall n.\ 0 \le n < size\_of\_ftype(ft) \Rightarrow ct(v, n) = none \wedge ct'(v,n) = undefined \\
    \updatecm(cm_1, (v, 0), ft, [\underbrace{\svundefined, \ldots, \svundefined}_{m}]) = cm_2
    \end{array}
   }
    \quad
    \\[1ex]
    \infer[node]
    {\seqq{ cm_1 \xrightarrow[tmap]{\mbox{node } x = e} cm_2}}
    {
    \begin{array}{c}
   id(x) = v \wedge \typeofexp(e, tmap)= ft \wedge \sizeofftype(ft) = sz \wedge \listgexp(e, ft) = ls \\
    \forall n.\ 0 \le n < sz  \Rightarrow  cm_1(v, n) = \svnone \\
    \updatecm(cm_1, (v, 0), ft, ls) = cm_2
    \end{array}    
    }
    \quad
    \\[1ex]
    \infer[register]
    {\seqq{ cm_1 \xrightarrow[tmap]{\mbox{reg } x : t\ e\ \mbox{with: reset=>}  (rsig, rval)} cm_2}}
    {
    \begin{array}{c}
   id(x) = v \wedge tmap(v)= (ft, ori) \wedge \sizeofftype(ft) = sz \wedge  rft = \typeofexp(rval, tmap) \\
   \typecompatible(ft, rft, false) \wedge \forall n.\ 0 \le n < sz  \Rightarrow  cm_1(v, n) = \svnone \\
   \updatecm(cm_1, (v, 0), ft, \listgexp(x, ft)) = cm_2
    \end{array}    
    }
    \quad
    \\[1ex]
    \infer[connect: exp]
    {\seqq{ cm_1 \xrightarrow[tmap]{r <= e} cm_2}}
    {
    \begin{array}{c}
    id(r)=(v, m) \wedge \typeofref(r, tmap)= (rft, ori) \ \wedge \\
     (ori = \duplex \vee ori = \sink) \wedge \sizeofftype(rft) = sz \wedge \typeofexp(e, tmap)= eft \\
    \typecompatible(rft, eft, false) \wedge \forall n.\ 0 \le n < sz \Rightarrow cm_1(v, m+n) \neq \svnone \\
%    \forall n.\ n \ge size\_of\_ftype(ft) \Rightarrow vm(v, n) =  vm'(v,n)\\
%    \forall n.\ 0 \le n < size\_of\_ftype(ft) \Rightarrow ct(v, n) = none \wedge ct'(v,n) = undefined \\
     \updatecm(cm_1, (v, m), rft, \listgexp(e, eft)) = cm_2
    \end{array}    
    }
    \quad
    \\[1ex]
    \infer[connect: ref]
    {\seqq{ cm_1 \xrightarrow[tmap]{r <= r'} cm_2}}
    {
    \begin{array}{c}
     id(r)=(v, m) \wedge  id(r')=(v', m') \wedge \typeofref(r, tmap)= (ft, ori) \ \wedge \\
     \sizeofftype(ft) = sz \wedge \typeofref(r', tmap)= (ft', ori')\ \wedge \\
     \listgexp(r, ft) = ls \wedge  \listgexp(r', ft') = ls'\ \wedge  \\
    (ori = \duplex \vee ori = \sink) \wedge (ori' = \duplex \vee ori' = \source) \\
    \typecompatible(ft, ft', false) \\
%    \forall n.\ n \ge size\_of\_ftype(ft) \Rightarrow vm(v, n) =  vm'(v,n)\\
%    \forall n.\ 0 \le n < size\_of\_ftype(ft) \Rightarrow ct(v, n) = none \wedge ct'(v,n) = undefined \\
	\forall n.\ 0 \le n < sz \Rightarrow (cm_1(v, m+n) \neq \svnone \wedge cm_1(v', m'+n) \neq \svnone) \\
     \updatecmref(cm_1, (v, m), ft, ls, (v', m'), ft', ls') = cm_2
%    \forall v'', n''.\ (v'' \neq v \vee n'' < m \vee n'' \ge m+sz) \wedge (v'' \neq v' \vee n'' < m' \vee n'' \ge m'+sz)  \Rightarrow \\
%    ct'(v'', n'') = ct(v'', n'') 
    \end{array}    
    }
    \\[1ex]
    \infer[invalid]
    {\seqq{cm_1 \xrightarrow[tmap]{r \mbox{ is invalid}} cm_2}}
    {
    \begin{array}{c}
    id(r) = (v, m) \wedge \typeofref(r, tmap) = (ft, ori)\\
    \invalidate(cm_1, (v,m), ft, ori) = cm_2
    \end{array}    
    }
    \\[1ex]
    \infer[when]
    {\seqq{cm_1 \xrightarrow[tmap]{\mbox{when } cond:\ sst \mbox{ else}:\ ssf } cm_2}}
    {
    \begin{array}{c}
    \typeofexp(cond, tmap) = (\gtype, (\uint, 1))\\
    cm \xrightarrow[tmap]{sst} cmt \wedge cm \xrightarrow[tmap]{ssf} cmf \wedge
    cm_2 = \mux(cond, cmt, cmf)
    \end{array}    
    }
  \end{gather*}
%  \vspace{-4mm}
  \caption{The relation $cm_1 \xrightarrow[tmap]{s} cm_2$}
  \label{fig:firrtl-sem}
\end{figure}
 
       
%Let $M$ be a {\hifirrtl} module and $G$ be a module graph. We define how $G = (vm, ct)$ \emph{matches} $M$ as follows: We define the relations $(vm, ct) \xrightarrow[tmap]{s} (vm', ct')$, where $s$ is a statement. Then $G$ matches $M$ if $(vm_{emp}, ct_{emp}) \xrightarrow[tmap]{ss} (vm, ct)$, where $ss$ is the statement sequence of $M$. 

%%%%% vm_match

\begin{figure}[htbp]
  \small
  \begin{gather*}
    \infer[\updatecm]
    {\updatecm(cm_1, id, ft, ls) = cm_2}
    {
    \begin{array}{c}
    id = (v, m) \wedge \sizeofftype(ft) = sz \\
    \forall n.\ 0 \le n < sz \Rightarrow cm_2(v, m+n) = ls[n] \\
    \forall v', n'.\ (v' \neq v \vee n' < m \vee n' \ge m+sz) \Rightarrow cm_2(v', n') = cm_1(v', n') \\
    \end{array}
    }
    \\[1ex]
    \infer[\updatecmref]
    {\updatecmref(cm_1, id, ft, ls, id', ft', ls') = cm_2}
    {
    \begin{array}{c}
    id = (v, m) \wedge id' = (v', m') \wedge \sizeofftype(ft) = sz\ \wedge \\
    \listgtyperef(ft', false) = lgt \\
    \forall n.\ (0 \le n < sz \wedge lgt[n] = (\_, false)) \Rightarrow \\
    \big(
%    \begin{array}{l}
    cm_2(v, m+n) = ls'[n] \wedge ((v',m') \neq (v, m)  \rightarrow cm_2(v', m' + n) = cm_1(v', m'+n)) 
%    \end{array}
    \big)
    \\ 
    \forall n.\ (0 \le n < sz \wedge lgt[n] = (\_, true)) \Rightarrow \\
    \big(
%    \begin{array}{l}
    cm_2(v', m'+n) = ls[n] \wedge  
    ((v, m) \neq (v', m')  \rightarrow cm_2(v, m + n) = cm_1(v, m+n))
%    \end{array}
    \big)
     \\
     \forall v'', n''.\ 
     \left(
     \begin{array}{l}
     (v'' \neq v \vee n'' < m \vee n'' \ge m+sz)\ \wedge \\
     (v'' \neq v' \vee n'' < m' \vee n'' \ge m'+sz)
     \end{array}
     \right)
       \Rightarrow cm_2(v'', n'') = cm_1(v'', n'') 
    \end{array}
    }
    \\[1ex]
    \infer[\invalidate]
    {\invalidate(cm_1, id, ft, ori) = cm_2}
    {
    \begin{array}{c}
    id = (v, m) \wedge \sizeofftype(ft) = sz \wedge \listgtyperef(ft, false) = lgt \\
    \forall n.\ 0 \le n < sz \wedge 
    \left(
    \begin{array}{l}
    (lgt[n] = (\_, false) \wedge (ori = \sink \vee ori = \duplex))\ \vee \\
    (lgt[n] = (\_, true) \wedge (ori = \source \vee ori = \duplex))
     \end{array}
    \right) \Rightarrow\\
    \hspace{1cm}  cm_2(v, m+n) = \svinvalid  \\
    \forall n.\ 0 \le n < sz \wedge 
    \left(
    \begin{array}{l}
    (lgt[n] = (\_, false) \wedge ori = \source)\ \vee \\
     (lgt[n] = (\_, true) \wedge ori = \sink)
     \end{array}
    \right)  
    \Rightarrow cm_2(v, m+n) = cm_1(v, m+n)  \\
    \forall v', n'.\ (v' \neq v \vee n' < m  \vee n' \ge m+sz)  \rightarrow cm_2(v', n') = cm_1(v', n')
    \end{array}
    }
  \end{gather*}
  \caption{$\updatecm$, $\updatecmref$, and $\invalidate$}
  \label{tab:update-vm-ct}
\end{figure}

\begin{figure}[htbp]
  \small
  \begin{gather*}
    \infer[\listgexp: ground]
    {\listgexp(e, ft) = l}
    {
    \begin{array}{c}
    	ft = (ground, t) \wedge l = [e]
    \end{array}
    }
    \\[1ex]
    \infer[\listgexp: mux]
    {list\_gexp(e, ft) = l}
    {
    \begin{array}{c}
    	e = mux(c, e1, e2) \wedge list\_gexp(e1, ft) = l1 \wedge list\_gexp(e1, ft) = l2\  \wedge\\
	size(l1) = size(l2) \wedge l = mux(c, l1, l2)
    \end{array}
    }
    \\[1ex]
    \infer[\listgexp: vector]
    {list\_gexp(e, ft) = l}
    {
    \begin{array}{c}
    	 e= r[n] \wedge type\_of\_ref(r) = (ft', ori)\ \wedge \\
	  size\_of\_ftype(ft) = sz \wedge l = list\_gexp(r, ft')[sz* n, sz*(n+1)-1]
    \end{array}
    }
    \\[1ex]
    \infer[list\_gexp: bundle]
    {list\_gexp(e, ft) = l}
    {
    \begin{array}{c}
    	 ft = (bundle, ff) \wedge l = list\_gexp\_fields(e, ff)
    \end{array}
    }
    \\[1ex]
    \infer[list\_gexp\_fields: nil]
    {list\_gexp\_fields(e, ff) = l}
    {
    \begin{array}{c}
    	 ff = nil \wedge l = []
    \end{array}
    }
   \\[1ex]
    \infer[list\_gexp\_fields: $\neq$nil]
    {list\_gexp\_fields(e, ff) = l}
    {
    \begin{array}{c}
    	 ff = (f, b, t, ff1) \wedge l = list\_gexp(e.f, t)\ ++\ list\_gexp\_fields(e, ff1)
    \end{array}
    }
  \end{gather*}
  \caption{$list\_gexp$ and $list\_gexp\_fields$}
  \label{fig:list-gexp}
\end{figure}


\begin{figure}[htbp]
  \small
  \begin{gather*}
    \infer[\listgtype: ground]
    {\seqq{list\_gtype(ft) = [gt]}}
    {\seqq{ft = (ground, gt)}}
    \quad
    \\[1ex]
    \inferii[list\_gtype: vector]
    {\seqq{list\_gtype(ft) = \underbrace{l'\ ++\ \cdots\ ++\ l'}_{m}}}
    {\seqq{ft = (vector, ft'[m])}}    {\seqq{l' = list\_gtype(ft')}}
    \quad
    \\[1ex]
    \inferii[list\_gtype: bundle]
    {\seqq{list\_gtype(ft) = l'}}
    {\seqq{ft = (bundle, ff)}}    {\seqq{l' = list\_gtype\_fields(ff)}}
    \quad
    \\[1ex]    
    \infer[list\_gtype\_fields: nil]
    {\seqq{list\_gtype\_fields(ff) = []}}
    {\seqq{f f = nil}}
    \quad
    \\[1ex]
    \inferiii[list\_gtype\_fields: $\neq$ nil]
    {\seqq{list\_gtype\_fields(ff) = l'\ {++}\ l''}}
    {\seqq{f f = (\_, \_, ft', ff')}} 
    {\seqq{list\_gtype(ft') = l'}} {\seqq{list\_gtype\_fields(ff') = l''}}  
    \quad
    \\[1ex]
    \infer[list\_gtype\_ref: ground]
    {\seqq{list\_gtype\_ref(ft, flipped) = [(gt, flipped)]}}
    {\seqq{ft = (ground, gt)}}
    \quad
    \\[1ex]
    \inferii[list\_gtype\_ref: vector]
    {\seqq{list\_gtype\_ref(ft, flipped) = \underbrace{l'\ ++\ \cdots\ ++\ l'}_{m}}}
    {\seqq{ft = (vector, ft'[m])}}    {\seqq{l' = list\_gtype\_ref(ft', flipped)}}
    \quad
    \\[1ex]
    \inferii[list\_gtype\_ref: bundle]
    {\seqq{list\_gtype\_ref(ft, flipped) = l'}}
    {\seqq{ft = (bundle, ff)}}    {\seqq{l' = list\_gtype\_ref\_fields(ff, flipped)}}
    \quad
    \\[1ex]    
    \infer[list\_gtype\_ref\_fields: nil]
    {\seqq{list\_gtype\_ref\_fields(ff, flipped) = []}}
    {\seqq{f f = nil}}
    \quad
    \\[1ex]
    \inferii[list\_gtype\_ref\_fields: $\neq$nil-1]
    {\seqq{list\_gtype\_ref\_fields(ff, flipped) = l'\ {++}\ l''}}
    {
    \begin{array}{c}
    	f f = (\_, flipped', ft', ff')\ \wedge\\
	 flipped' = false
    \end{array}
    }
    {
    \begin{array}{c}
    	list\_gtype\_ref(ft', flipped) = l'\ \wedge \\
    	list\_gtype\_ref\_fields(ff', flipped) = l''
    \end{array}
    } 
    \quad
    \\[1ex]
    \inferii[list\_gtype\_ref\_fields: $\neq$nil-2]
    {\seqq{list\_gtype\_ref\_fields(ff, flipped) = l'\ {++}\ l''}}
    {
    \begin{array}{c}
    	f f = (\_, flipped', ft', ff')\ \wedge\\
	 flipped' = true
    \end{array}
    }
    {
    \begin{array}{c}
    	list\_gtype\_ref(ft', \neg flipped) = l'\ \wedge \\
    	list\_gtype\_ref\_fields(ff', flipped) = l''
    \end{array}
    } 
  \end{gather*}
%  \vspace{-4mm}
  \caption{$list\_gtype$, $list\_gtype\_fields$, $list\_gtype\_ref$, and $list\_gtype\_ref\_fields$}
  \label{fig:size-list-ftype-fields}
\end{figure}


\begin{figure}[htbp]
  \small
  \begin{gather*}
    \infer[type\_compatible: exp-gt]
    {type\_compatible(ft, ft', flipped)}
    {
    \begin{array}{c}
    	ft = (ground, t) \wedge ft' = (ground, t) \\
	 t \in \{UInt, SInt\} \times \Nat \cup \{Clock, Reset, AsyncReset\}
    \end{array}
    }
    \\[1ex]
    \infer[type\_compatible: imp-gt-1]
    {type\_compatible(ft, ft', flipped)}
    {
    \begin{array}{c}
    	ft = (ground, t) \wedge ft' = (ground, t') \\
	(t = (UIntImp, n) \wedge t' = (UInt, n')) \vee (t = (SIntImp, n) \wedge t' = (SInt, n')) \\
	flipped = false \rightarrow n \ge n'
    \end{array}
    }
    \\[1ex]
    \infer[type\_compatible: imp-gt-2]
    {type\_compatible(ft, ft', flipped)}
    {
    \begin{array}{c}
    	ft = (ground, t) \wedge ft' = (ground, t') \\
	(t = (UInt, n) \wedge t' = (UIntImp, n')) \vee (t = (SInt, n) \wedge t' = (SIntImp, n')) \\
	flipped = true \rightarrow n \le n'
    \end{array}
    }
    \\[1ex]
    \infer[type\_compatible: imp-gt-3]
    {type\_compatible(ft, ft', flipped)}
    {
    \begin{array}{c}
    	ft = (ground, t) \wedge ft' = (ground, t') \\
	(t = (UIntImp, n) \wedge t' = (UIntImp, n')) \vee (t = (SIntImp, n) \wedge t' = (SIntImp, n')) \\
	(flipped = false \wedge n \ge n') \vee (flipped = true \wedge n \le n')
    \end{array}
    }
    \\[1ex]
    \infer[type\_compatible: vector]
    {type\_compatible(ft, ft', flipped)}
    {
    \begin{array}{c}
    	ft = (vector, t[n]) \wedge ft' = (vector, t'[n]) \\
	type\_compatible(t, t', flipped)
    \end{array}
    }
    \\[1ex]
    \infer[type\_compatible: bundle]
    {type\_compatible(ft, ft', flipped)}
    {
    \begin{array}{c}
    	ft = (bundle, ff) \wedge ft' = (bundle, ff') \\
	type\_compatible\_fields(ff, ff', flipped)
    \end{array}
    }
    \\[1ex]
    \infer[type\_compatible\_fields: nil]
    {type\_compatible\_fields(ff, ff', flipped)}
    {
    \begin{array}{c}
    	ff =nil \wedge ff' = nil
    \end{array}
    }
    \\[1ex]
    \infer[type\_compatible\_fields: $\neq$-nil-false]
    {type\_compatible\_fields(ff, ff', flipped)}
    {
    \begin{array}{c}
    	ff = (f, b, t, ff1) \wedge ff' = (f', b, t', ff1') \wedge b = false \\
	type\_compatible(t, t', flipped) \\
	type\_compatible\_fields(ff1, ff1', flipped)
    \end{array}
    }
    \\[1ex]
    \infer[type\_compatible\_fields: $\neq$-nil-true]
    {type\_compatible\_fields(ff, ff', flipped)}
    {
    \begin{array}{c}
    	ff = (f, b, t, ff1) \wedge ff' = (f', b, t', ff1') \wedge b = true \\
	type\_compatible(t, t', \neg flipped) \\
	type\_compatible\_fields(ff1, ff1', flipped)
    \end{array}
    }
  \end{gather*}
  \caption{$type\_compatible$ and $type\_compatible\_fields$}
  \label{fig:type-compatible}
\end{figure}




\begin{figure}[htbp]
  \small
  \begin{gather*}
   \inferii[\typeofref: id]
    {\seqq{\typeofref(r, tmap) = (ft, ori)}}
    {\seqq{r = x}} {\seqq{tmap(id(x)) = (ft, ori)}}
    \quad
    \\[1ex]
   \inferiii[\typeofref: vector-1]
    {\seqq{\typeofref(r, tmap) = (t, ori)}}
    {\seqq{r = r'[n]}} {\seqq{\typeofref(r', tmap) =  (t[m], ori)}} {\seqq{n < m}}
    \quad
    \\[1ex]
   \inferiii[\typeofref: vector-2]
    {\seqq{\typeofref(r, tmap) = none}}
    {\seqq{r = r'[n]}} {\seqq{\typeofref(r', tmap) = (t[m], ori)}} {\seqq{n \ge m}}
    \quad
    \\[1ex]
   \infer[\typeofref: field]
    {\seqq{\typeofref(r, tmap) = (ft, ori)}}
    {
    \begin{array}{c}
    r = r'.f \\
    \typeofref(r', tmap) = (ft', ori') \\
    \typeoffield(f, ft', ori') = (ft, ori)
    \end{array}
    }
    \quad
%    \seqq{r = r'.f}} {\seqq{type\_of\_ref(r', tmap) = (ft', ori')}} {\seqq{type\_of\_field(f, ft', ori') = (ft, ori)}}
    \\[1ex]
   \infer[\typeoffield: \nil]
    {\seqq{\typeoffield(f, ff, ori) = none}}
    {\seqq{f f = \nil}}
    \quad
    \\[1ex]
   \infer[\typeoffield: neq]
    {\seqq{\typeoffield(f, ff, ori) = (ft, ori')}}
    {
    \begin{array}{c}
    f f = (f', flipped, t, ff')\wedge f \neq f'\\
    \typeoffield(f, ff', ori) = (ft, ori')
    \end{array}
    }
    \quad
    \\[1ex]
   \infer[\typeoffield: eq-1]
    {\seqq{\typeoffield(f, ff, ori) = (t, ori)}}
    {
    \begin{array}{c}
    f f = (f', flipped, t, ff') \wedge f = f' \\
    flipped = false \vee (flipped = true \wedge ori = \duplex)
    \end{array}
    }
    \quad
    \\[1ex]
   \infer[\typeoffield: eq-2]
    {\seqq{\typeoffield(f, ff, ori) = (t, \sink)}}
    {
    \begin{array}{c}
    f f = (f', flipped, t, ff') \wedge f = f' \\
    flipped = true \wedge ori = \source
    \end{array}
    }
    \quad
    \\[1ex]
   \infer[\typeoffield: eq-3]
    {\seqq{\typeoffield(f, ff, ori) = (t, \source)}}
    {
    \begin{array}{c}
    f f = (f', flipped, t, ff') \wedge f = f' \\
    flipped = true \wedge ori = \sink
    \end{array}
    }
  \end{gather*}
%  \vspace{-4mm}
  \caption{The procedures $\typeofref$ and $\typeoffield$}
  \label{fig:type-ref}
\end{figure}


